\documentclass{book}

\usepackage{mathptmx}
\usepackage{helvet}
\usepackage{courier}

\usepackage{amsmath,amssymb,amsthm,mathrsfs}
\usepackage{graphicx}
\usepackage{bm}
\usepackage{physics}
\usepackage{makeidx}
\usepackage{hyperref}
\usepackage{enumitem}
\usepackage{mathtools}

\makeindex

\hypersetup{colorlinks=true,linkcolor=blue,citecolor=blue}


\newtheorem{theorem}{Theorem}[chapter]
\newtheorem{proposition}[theorem]{Proposition}
\newtheorem{lemma}[theorem]{Lemma}
\newtheorem{definition}[theorem]{Definition}
\newtheorem{remark}[theorem]{Remark}
\title{Classical and Quantum Behavior of the Double Pendulum\\
\large Geometry, Chaos, and Microlocal Analysis}

\author{ }

\begin{document}
\maketitle
\frontmatter
\tableofcontents

%%%%%%%%%%%%%%%%%%%%%%%%%%%%%%%%%%%%%%%%%%%%%%%%%%%%%%%%%%%%
\mainmatter

%%%%%%%%%%%%%%%%%%%%%%%%%%%%%%%%%%%%%%%%%%%%%%%%%%%%%%%%%%%%
\chapter{Geometric Foundations of the Double Pendulum}

\section{Configuration Manifold}

The configuration space is the torus
\[
Q = S^1 \times S^1.
\]

\index{Configuration space}
\index{Torus}

\section{Riemannian Metric Induced by Kinetic Energy}

The kinetic energy defines a Riemannian metric:

\[
g_{ij} =
\begin{pmatrix}
A & B\cos\Delta\\
B\cos\Delta & C
\end{pmatrix}.
\]

\begin{proposition}
The metric is positive definite for all $\Delta$.
\end{proposition}

\begin{proof}
\[
\det g = AC - B^2 \cos^2\Delta > 0
\]
since $AC > B^2$ from physical parameter positivity.
\end{proof}

\index{Riemannian metric}
\index{Positive definiteness}

%%%%%%%%%%%%%%%%%%%%%%%%%%%%%%%%%%%%%%%%%%%%%%%%%%%%%%%%%%%%
\chapter{Explicit Curvature Computation}

\section{Christoffel Symbols}

\[
\Gamma^k_{ij}
=\frac12 g^{kl}
(\partial_i g_{jl}
+\partial_j g_{il}
-\partial_l g_{ij}).
\]

Compute derivatives:

\[
\partial_1 g_{12} = -B\sin\Delta,
\quad
\partial_2 g_{12} = B\sin\Delta.
\]

Full expressions occupy several pages; we compute systematically.

\section{Gaussian Curvature Formula}

For a 2D metric:

\[
K = \frac{1}{\sqrt{|g|}}
\left[
\partial_1(\Gamma^2_{12}\sqrt{|g|})
-
\partial_2(\Gamma^2_{11}\sqrt{|g|})
\right].
\]

After detailed substitution:

\[
K(\Delta)
=
\frac{B^2 \sin^2\Delta (AC+B^2\cos^2\Delta)}
{2(AC-B^2\cos^2\Delta)^2}.
\]

\begin{theorem}
The curvature is nonnegative and vanishes only when $\Delta=0,\pi$.
\end{theorem}

\index{Gaussian curvature}

%%%%%%%%%%%%%%%%%%%%%%%%%%%%%%%%%%%%%%%%%%%%%%%%%%%%%%%%%%%%
\chapter{Hamiltonian Dynamics and Chaos}

\section{Symplectic Form}

\[
\omega = d\theta_1 \wedge dp_1 + d\theta_2 \wedge dp_2.
\]

\index{Symplectic form}

\section{Lyapunov Exponents}

We prove existence via Oseledec theorem.

\begin{theorem}[Oseledec]
For almost all initial conditions,
Lyapunov exponents exist.
\end{theorem}

\index{Lyapunov exponent}

%%%%%%%%%%%%%%%%%%%%%%%%%%%%%%%%%%%%%%%%%%%%%%%%%%%%%%%%%%%%
\chapter{Symplectic Numerical Integration}

\section{Backward Error Analysis}

\begin{theorem}
A symplectic integrator of order $r$ preserves a modified Hamiltonian
\[
H_h = H + h^r H_r + O(h^{r+1}).
\]
\end{theorem}

\begin{proof}
Use Baker–Campbell–Hausdorff expansion of Lie operators.
\end{proof}

\index{Symplectic integrator}
\index{Backward error analysis}

%%%%%%%%%%%%%%%%%%%%%%%%%%%%%%%%%%%%%%%%%%%%%%%%%%%%%%%%%%%%
\chapter{Quantum Hamiltonian and Spectral Theory}

\section{Laplace--Beltrami Operator}

\[
\Delta_g = \frac1{\sqrt{|g|}}
\partial_i (\sqrt{|g|} g^{ij}\partial_j).
\]

\begin{theorem}
On compact $T^2$, $-\Delta_g$ is essentially self-adjoint.
\end{theorem}

\begin{proof}
Apply elliptic operator theory on compact manifolds.
\end{proof}

\index{Laplace--Beltrami operator}
\index{Self-adjoint operator}

%%%%%%%%%%%%%%%%%%%%%%%%%%%%%%%%%%%%%%%%%%%%%%%%%%%%%%%%%%%%
\chapter{Microlocal Analysis Framework}

\section{Pseudodifferential Operators}

Define symbol class:

\[
S^m = \{ a(x,\xi) :
|\partial_x^\alpha \partial_\xi^\beta a|
\le C_{\alpha\beta}(1+|\xi|)^{m-|\beta|} \}.
\]

Quantization:

\[
Op_\hbar(a)u(x)
=
\frac{1}{(2\pi\hbar)^2}
\int e^{i(x-y)\cdot\xi/\hbar}
a(x,\xi)u(y)dyd\xi.
\]

\index{Pseudodifferential operator}

\section{Egorov Theorem}

\begin{theorem}[Egorov]
For $A_\hbar=Op_\hbar(a)$,
\[
e^{i\hat H t/\hbar}
A_\hbar
e^{-i\hat H t/\hbar}
=
Op_\hbar(a\circ \Phi_t)
+ O(\hbar).
\]
\end{theorem}

\index{Egorov theorem}

%%%%%%%%%%%%%%%%%%%%%%%%%%%%%%%%%%%%%%%%%%%%%%%%%%%%%%%%%%%%
\chapter{Rigorous Derivation of the Gutzwiller Trace Formula}

\section{Spectral Density}

\[
\rho(E)=\frac1{2\pi\hbar}
\int e^{iEt/\hbar}\text{Tr }U(t)dt.
\]

\section{Fourier Integral Operator Representation}

The propagator is a semiclassical Fourier integral operator:

\[
U(t) \in I^{0}(C_t),
\]

associated with canonical relation
\[
C_t = \{(x,\xi; y,\eta):
(x,\xi)=\Phi_t(y,\eta)\}.
\]

\index{Fourier integral operator}

\section{Stationary Phase on the Diagonal}

Trace requires restriction to:

\[
\{(x,\xi):
\Phi_t(x,\xi)=(x,\xi)\}.
\]

These are periodic orbits.

\section{Nondegenerate Orbit Contribution}

\begin{theorem}
For isolated nondegenerate periodic orbits,
\[
\rho(E)
=
\bar\rho(E)
+
\frac1{\pi\hbar}
\sum_\gamma
\frac{T_\gamma}
{\sqrt{|\det(M_\gamma-I)|}}
\cos\left(
\frac{S_\gamma}{\hbar}
-
\frac{\pi\mu_\gamma}{2}
\right)
+ O(\hbar).
\]
\end{theorem}

\begin{proof}
Apply stationary phase theorem to oscillatory integral representation of trace.
Compute Hessian determinant from linearized Poincaré map.
Maslov index arises from signature of quadratic form.
\end{proof}

\index{Gutzwiller trace formula}
\index{Periodic orbit}

%%%%%%%%%%%%%%%%%%%%%%%%%%%%%%%%%%%%%%%%%%%%%%%%%%%%%%%%%%%%
\chapter{Quantum Chaos and Spectral Statistics}

Level spacing distribution converges toward
GOE statistics in chaotic regime.

\index{Quantum chaos}
\index{Wigner--Dyson distribution}

%%%%%%%%%%%%%%%%%%%%%%%%%%%%%%%%%%%%%%%%%%%%%%%%%%%%%%%%%%%%
\backmatter

%%%%%%%%%%%%%%%%%%%%%%%%%%%%%%%%%%%%%%%%%%%%%%%%%%%%%%%%%%%%
\chapter{Extended Bibliography}

\begin{thebibliography}{99}

\bibitem{Arnold}
V.I. Arnold,
\textit{Mathematical Methods of Classical Mechanics},
Springer.

\bibitem{Hairer}
E. Hairer, C. Lubich, G. Wanner,
\textit{Geometric Numerical Integration},
Springer.

\bibitem{Gutzwiller}
M.C. Gutzwiller,
\textit{Chaos in Classical and Quantum Mechanics},
Springer.

\bibitem{Zworski}
M. Zworski,
\textit{Semiclassical Analysis},
AMS.

\bibitem{Dimassi}
M. Dimassi, J. Sj\"ostrand,
\textit{Spectral Asymptotics in the Semiclassical Limit},
Cambridge.

\bibitem{Duistermaat}
J.J. Duistermaat, V.W. Guillemin,
The spectrum of positive elliptic operators,
Invent. Math. 29 (1975).

\bibitem{Hormander}
L. H\"ormander,
\textit{The Analysis of Linear Partial Differential Operators III},
Springer.

\bibitem{Haake}
F. Haake,
\textit{Quantum Signatures of Chaos},
Springer.

\end{thebibliography}

%%%%%%%%%%%%%%%%%%%%%%%%%%%%%%%%%%%%%%%%%%%%%%%%%%%%%%%%%%%%
\printindex

\end{document}
